\documentstyle[%


righttag%


,11pt%


,amssymb%


,russian%               remove in Eng.


,izvru%                 Set izven% in Eng.


]{amsart}





% {>}





\newcommand{\tts}[1]{}


\newcommand{\Mod}{\operatorname{mod}}


\newcommand{\mes}{\operatorname{mes}}


\renewcommand{\baselinestretch}{0.98}





\theoremstyle{plain}


\theorembodyfont{\it}


\newtheorem{thms}{\hskip\parindent Теорема}[section]





\theoremstyle{definition}


\theorembodyfont{\rm}


\newtheorem{defns}{\hskip\parindent Определение}[section]





%\hoffset=-15mm%                  For Eng.


\setcounter{page}{49}


\god{2005}


\nomer{4~(515)} %                        Remove in Eng.


%\nomer{Vol.\,49, No.\,4}%               Set in Eng.


%\str{\pageref{begin}--\pageref{end}}%  Set in Eng.


\udk{517.956}





\author{\it В.Н.\,Павленко, Д.К.\,Потапов}


% NB:   ^^ remove in Eng.


\title{Аппроксимация краевых задач эллиптического типа


со спектральным параметром и разрывной нелинейностью}





\date{}


%\pagestyle{myheadings}%         Set in Eng.


\pagestyle{plain} %              Remove in Eng.


\begin{document}


%\thispagestyle{plain}%          Set in Eng.


\maketitle


\label{begin}





\setcounter{section}{1}











{\bf Введение.}


Разрывные нелинейности в интегральных и дифференциальных уравнениях часто


возникают как идеализации непрерывных нелинейностей, имеющих участки быстрого


роста по фазовой переменной. При этом удобно считать, что непрерывные


нелинейности зависят от малого параметра $\varepsilon$ и в пределе при


$\varepsilon\to 0$ получается разрывная нелинейность. Естественным является


вопрос о близости решений аппроксимирующего уравнения и предельной задачи. На


необходимость исследований в этом направлении указано в~\cite{krasn2}. Данная


проблема для эллиптических краевых задач с разрывными нелинейностями изучалась


в~\cite{pavl4} в предположении, что нелинейность, входящая в уравнение,


ограничена, а дифференциальная часть вместе с граничным условием порождает


коэрцитивный оператор.





В данной работе рассматриваются резонансные эллиптические краевые задачи со


спектральным параметром и ограниченной разрывной нелинейностью. Резонансность


краевой задачи означает, что ядро дифференциального оператора с соответствующим


краевым условием ненулевое.


В~\cite{pot1} вариационным методом получены предложения о существовании луча


положительных собственных значений для таких задач. Аппроксимирующая задача


получается из исходной малыми возмущениями спектрального параметра и


непрерывными по фазовой переменной аппроксимациями разрывной нелинейности. Цель


данной работы -- установить при определенных условиях сходимость решений


аппроксимирующих задач к решениям исходной.


\medskip








{\bf 1. Постановка задачи.}


Пусть $\Omega$ --- ограниченная область в ${\bold R}^n$ c границей $\Gamma$ класса


${\bold C}_{2,\alpha}$, $0<\alpha \leq 1$ (\cite{lad}, с.\,23), гиперповерхности


$$


S_i=\{(x, u)\in {\bold R}^{n+1} : u=\varphi_i(x),\ x\in\overline{\Omega}\},


$$


$\varphi_i\in {\bold W}_q^2(\Omega)$, $q>n$, $i=\overline{1, m}$, попарно


не пересекаются. Для определенности считается, что $\varphi_i(x)<\varphi_{i+1}(x)$


для любого $x\in\overline{\Omega}$ и $i=\overline{1, m-1}$. Так как $q>n$, то


$\varphi_i(x)\in {\bold C}_1(\overline{\Omega})$ (\cite{sob}, с.\,74). 


Отсюда и из последнего


неравенства следует существование $d>0$ такого, что для любого


$x\in\overline{\Omega}$ отрезки $[\varphi_i(x)-d,\; \varphi_i(x)+d]$,


$i=\overline{1, m}$, попарно не пересекаются.





Поверхности $S_i$, $i=\overline{1, m}$, разбивают область


$D=\Omega\times {\bold R}$ на непересекающиеся подобласти


\begin{align*}


D_0&=\{(x, u)\in D : u<\varphi_1(x)\},\\


D_i&=\{(x, u)\in D : \varphi_i(x)<u<\varphi_{i+1}(x)\},\ i=\overline{1, m-1},\\


D_m&=\{(x, u)\in D : u>\varphi_m(x)\}.


\end{align*}





На $\overline{D_i}$ заданы каратеодориевы функции (\cite{kar}, с.\,148)


$g_i(x, u)$ такие, что для почти всех $x\in\Omega$


\begin{equation}


\label{odin}


|g_i(x, u)|\leq a(x) \ \; \forall u,\ \ (x,u)\in\overline{D_i},


\end{equation}


где $a\in {\bold L}_q(\Omega)$, $q>n$, $i=\overline{0, m}$.





Функция $g: \overline{D}\rightarrow {\bold R}$ суперпозиционно измерима


(\cite{kar}, с.\,149)


и на $D_i$ совпадает с $g_i(x, u)$, причем для почти всех $x\in\Omega$, если


$(x, u)\in S_i$, то $g(x, u)$ принадлежит отрезку с концами $g_{i-1}(x, u)$ и


$g_i(x, u)$. Будем предполагать, что для почти всех $x\in\Omega$


верно равенство


\begin{equation}


\label{dva}


g(x, 0)=0


\end{equation}


и


\begin{equation}


\label{n3}


g_i(x, u)\geq g_{i-1}(x, u),


\end{equation}


если $(x,u)\in S_i$.





Фиксируем последовательность положительных чисел $(\delta_k)$, сходящуюся к


нулю и ограниченную сверху определенным выше числом $d$.





Нелинейность $g(x, u)$ аппроксимируется последовательностью каратеодориевых


функций $(g^k(x, u))$ таких, что для почти всех $x\in\Omega$


\begin{itemize}


\item[(i)] $g^k(x, u)=g(x, u)$, если $|\varphi_i(x)-u|>\delta_k$ для любого


$i=\overline{1, m}$;


\item[(ii)]


\begin{equation}


\label{tri}


|g^k(x, u)|\leq a(x) \ \; \forall u\in {\bold R},


\end{equation}


где функция $a$ из оценки (\ref{odin}).


\end{itemize}





Исходная краевая задача с параметром $\lambda$ имеет вид


\begin{gather}


\label{uravn}


Lu(x)=\lambda g(x, u(x)), \ \ x\in\Omega,\\


\label{gran}


Bu|_{ \Gamma}=0,


\end{gather}


где


$Lu(x)\equiv-\sum\limits_{i,j=1}^n {(a_{ij}(x)u_{x_i})}_{x_j}+


c(x)u(x)$ ---  равномерно эллиптический


формально самосопряженный дифференциальный оператор в области


$\Omega$ с коэффициентами


$a_{ij} \in {\bold C}_{1,\alpha}(\overline{\Omega})$,


$c \in {\bold C}_{0,\alpha}(\overline{\Omega})$, $0<\alpha\leq 1$


(\cite{lad}, с.\,21).


При этом (\ref{gran}) --- одно из следующих основных граничных условий:


Дирихле, если $Bu\equiv u$;


Неймана, если


$Bu=\frac{\partial u }{\partial n_L}\equiv \sum\limits_{i,j=1}^n


a_{ij}(x)u_{x_i}\cos(n,x_j)$,


$n$ --- внешняя нормаль к границе $\Gamma$,


$\cos(n,x_j)$ --- направляющие косинусы нормали $n$;


третье краевое, если


$Bu=\frac{\partial u }{\partial n_L}+\sigma (x)u(x)$ c


$\sigma \in{\bold C}_{1,\alpha}(\Gamma)$, неотрицательной и не равной


тождественно нулю на $\Gamma$.





\begin{defns}


{\it Сильным решением} задачи (\ref{uravn})--(\ref{gran})


называется функция $u \in{\bold W}_s^2(\Omega)$, $s{>}1$, которая


удовлетворяет для почти


всех $x\in\Omega$ уравнению (\ref{uravn}) и для которой след $Bu(x)$ на


$\Gamma$ равен нулю.


\end{defns}





\begin{defns}[{\series{m}\shape{n}\selectfont \cite{krasn3}}]


Сильное решение $u(x)$ задачи (\ref{uravn})--(\ref{gran})


называется {\it полуправильным}, если для почти всех $x \in \Omega$ значение


$u(x)$ является точкой непрерывности функции $g(x, \cdot)$.


\end{defns}





\begin{defns}


Число $\lambda$ называется {\it собственным значением} задачи


(\ref{uravn})--(\ref{gran}), если существует ненулевое сильное решение $u(x)$


этой задачи. При этом $u(x)$ называют {\it собственной функцией},


соответствующей $\lambda$.


\end{defns}





\begin{defns}


Пусть $f:{\bold R} \to {\bold R}$. Назовем $u \in{\bold R}$


{\it прыгающим разрывом}


функции $f$, если $f(u-)<f(u+)$, где $f(u \pm)=\lim \limits_{s \to u \pm} f(s)$.


\end{defns}





В силу равенства (\ref{dva}) функция, почти всюду на $\Omega$ равная нулю,


является сильным решением задачи (\ref{uravn})--(\ref{gran}).





Заметим, что из неравенства (\ref{n3}) следует, что для почти всех $x\in\Omega$


точки разрыва функции $g(x, \cdot)$ прыгающие.





Пусть $X={\bold H}^1_0(\Omega)$, если (\ref{gran}) --- граничное условие


Дирихле, и $X={\bold H}^1(\Omega)$, если (\ref{gran}) --- граничное условие


Неймана или третье краевое условие.





Сопоставим краевой задаче (\ref{uravn})--(\ref{gran}) функционал $J^\lambda$,


определенный на $X$ следующим образом:


$$


J^{\lambda}(u)=J_1(u)-\lambda\int_{\Omega}dx


\int_0^{u(x)}g(x,s)ds,


$$


где


$$


J_1(u)=\frac{1}{2}\sum\limits^{n}_{i,j=1}


\int_{\Omega}a_{ij}(x)u_{x_i}u_{x_j}dx+\frac{1}{2}\int_{\Omega}c(x)u^2(x)dx


$$


в случае граничного условия Дирихле или Неймана;


$$


J_1(u)=\frac{1}{2}\sum\limits^{n}_{i,j=1}


\int_{\Omega}a_{ij}(x)u_{x_i}u_{x_j}dx+\frac{1}{2}\int_{\Omega}c(x)u^2(x)dx+


\frac{1}{2}\int_{\Gamma}\sigma(s)u^2(s)ds


$$


в случае третьего краевого условия.





Предполагается, что


\begin{itemize}


\item[(iii)] линейное пространство $N(L)$ решений задачи


\begin{gather}


\label{n6}


Lu=0,\quad x\in\Omega,\\


\label{n7}


Bu|_{\Gamma}=0


\end{gather}


одномерно и $\psi (x)$ --- базисная функция этого подпространства.


\end{itemize}





В дальнейшем используется следующий результат, который является частным


случаем теоремы 1.3 из~\cite{vadim}.





\begin{thms} Пусть выполнены условия


\begin{itemize}


\item[1)] $J_1(u)\geq 0$ $\forall u\in X;$


\item[2)] функция $g : \Omega\times {\bold R} \rightarrow {\bold R}$ борелева


$(\Mod 0)$ {\em (\cite{kar}, с. 166)}, для почти всех $x\in\Omega$ сечение


$g(x, \cdot)$ имеет


разрывы только первого рода, $g(x, u)\in [g_-(x, u), g_+(x, u)]$,


$g_-(x, u)=\underset{\eta\to u}\varliminf{}g(x, \eta)$,


$g_+(x, u)=\underset{\eta\to u}\varlimsup{}g(x, \eta)$ и верна оценка


$|g(x, u)|\leq a(x)$ $\forall u\in {\bold R}$, $a\in {\bold L}_q(\Omega)$, $q>n;$


\item[3)] линейное пространство $N(L)$ решений задачи $(\ref{n6})$--$(\ref{n7})$


ненулевое, $\lambda>0$ и


$$


\lim\limits_{u \in N(L),\;\| u \| \to +\infty}


\displaystyle\int_{\Omega}dx\int_{0}^{u(x)} g(x,s)ds=-\infty.


$$


\end{itemize}





Тогда существует $u_0\in X$, для которого


\begin{equation}


\label{n8}


J^{\lambda}(u_0)=\inf_{v\in X}J^{\lambda}(v)\equiv d_{\lambda},


\end{equation}


причем любое такое $u_0$ принадлежит ${\bold W}_q^2(\Omega)$


и удовлетворяет граничному условию $(\ref{gran})$.


Если дополнительно предположить, что для почти всех $x\in\Omega$ функция


$g(x, \cdot)$ имеет только прыгающие разрывы, то любое $u_0$, удовлетворяющее


$(\ref{n8})$, является полуправильным решением задачи 


$(\ref{uravn})$--$(\ref{gran})$.


\label{t3}


\end{thms}





В~\cite{pot1} показано, что если дополнительно к условиям 1)--3)


теоремы \ref{t3} $g(x, 0)=0$ для почти всех $x\in\Omega$ 


\begin{itemize}


\item[(iv)] существует $\wh{u} \in X$, для которого


$$


\int_{\Omega}dx\int_0^{\wh{u}(x)}g(x, s)ds>0,


$$


\end{itemize}


то найдется такое $\lambda_0>0$, что для любого $\lambda>\lambda_0$


число $d_{\lambda}$ в (\ref{n8}) отрицательно, поэтому $u_0$ в


равенстве (\ref{n8}) отлично от нулевого элемента $X$.





Будем предполагать, что условие (iv) выполнено. Кроме того, пусть для базисной


функции $\psi$ пространства $N(L)$ верны неравенства (условия Ландесмана--Лазера)


\begin{equation}


\int_{\psi<0}g_+(x)\psi(x)dx+\int_{\psi>0}g_-(x)\psi(x)dx>0>


\label{n9}


\int_{\psi>0}g_+(x)\psi(x)dx+\int_{\psi<0}g_-(x)\psi(x)dx,


\end{equation}


где $g_+(x)=\lim\limits_{u\to +\infty}g(x, u)$,


$g_-(x)=\lim\limits_{u\to -\infty}g(x, u)$ (предполагается, что для почти всех


$x\in\Omega$ данные пределы существуют). Как показано в~\cite{vadim}, условия


Ландесмана--Лазера влекут равенство


\begin{equation}


\label{n10}


\lim\limits_{u\in N(L),\; \|u\|\to +\infty}


\int_{\Omega}dx\int_0^{u(x)}g(x, s)ds=-\infty.


\end{equation}





Если дополнительно потребовать неотрицательность $J_1(u)$ на $X$, то из теоремы


\ref{t3} и замечания к ней при сделанных предположениях относительно задачи


(\ref{uravn})--(\ref{gran}) следует существование $\lambda_0>0$ такого, что


при $\lambda>\lambda_0$ \ $\inf\limits_{v\in X}J^\lambda(v)=d_{\lambda}<0$ и


найдется $u\in X$, для которого $J^\lambda(u)=d_\lambda$, причем любое такое


$u$ принадлежит ${\bold W}_q^2(\Omega)$ и является ненулевым полуправильным


решением задачи (\ref{uravn})--(\ref{gran}).





Зафиксируем $\lambda>\lambda_0$ и пусть числовая последовательность


$(\lambda_k)$, $\lambda_k>0$, сходится к $\lambda$.


Рассмотрим аппроксимирующую задачу


\begin{gather}


\label{uravn1}


Lu(x)=\lambda_k g^k(x,u(x)), \quad x\in\Omega,\\


\label{gran1}


Bu|_{ \Gamma}=0,


\end{gather}


где аппроксимирующая последовательность каратеодориевых 


функций $(g^k(x, u))$ определена выше.





Положим


$$


J_k(u)=J_1(u)-\lambda_k\int_{\Omega}dx\int_0^{u(x)}g^k(x,s)ds,


\quad {\bold N}\ni k\geq 2.


$$


По построению $g^k(x, u)$ каратеодориева и для нее при почти всех $x\in\Omega$


верна оценка (\ref{tri}) с функцией $a$ из оценки (\ref{odin}).





Покажем, что (\ref{n10}) для любого $k\geq 2$ влечет равенство


\begin{equation}


\label{n13}


\lim\limits_{u\in N(L),\; \|u\|\to +\infty}\int_{\Omega}dx


\int_0^{u(x)}g^k(x, s)ds=-\infty.


\end{equation}


Действительно, для произвольного $u\in X$


\begin{multline*}


\bigg|\int\limits_{\Omega}dx\int\limits_0^{u(x)}g^k(x, s)ds-


\int\limits_{\Omega}dx\int\limits_0^{u(x)}g(x, s)ds\bigg|\leq \\


%


\le \sum_{i=1}^m\int\limits_{\Omega}dx


\int\limits_{\varphi_i(x)-\delta_k}^{\varphi_i(x)+\delta_k}


|g^k(x, s)-g(x, s)|ds\leq \\


%


\le \sum_{i=1}^m\int\limits_{\Omega}dx


\int\limits_{\varphi_i(x)-\delta_k}^{\varphi_i(x)+\delta_k}2a(x)ds=


4m\delta_k\int_{\Omega}a(x)dx.


\end{multline*}


Данная оценка верна для любого $u\in N(L)$. Поэтому из (\ref{n10}) следует


(\ref{n13}), поскольку для любого $u\in N(L)$


\begin{multline*}


\int_{\Omega}dx\int_0^{u(x)}g^k(x, s)ds=


\int_{\Omega}dx\int_0^{u(x)}g(x, s)ds+


\int_{\Omega}dx\int_0^{u(x)}(g^k(x, s)-g(x, s))ds\leq \\


%


\le\int_{\Omega}dx\int_0^{u(x)}g(x, s)ds+


4m\delta_k\int_{\Omega}a(x)dx.


\end{multline*}





Из теоремы \ref{t3} следует, что для любого натурального $k\geq 2$ существует


такое $u_k\in X$, что $J_k(u_k)=\inf\limits_{v\in X}J_k(v)$ и $u_k$


является сильным решением задачи (\ref{uravn1})--(\ref{gran1}).





Рассматривается проблема о близости решений аппроксимирующей задачи $u_k$ к


решениям исходной задачи (\ref{uravn})--(\ref{gran}).


\addtocounter{section}{1}


\setcounter{thms}{0}


\medskip





{\bf 2. Формулировка и доказательство основного результата}





\begin{thms}


Предположим, что


\begin{itemize}


\item[1)] $J_1(u)\geq 0$ $\forall u\in X;$


\item[2)] выполнены оценка $(\ref{odin})$, равенство $(\ref{dva})$,


неравенство $(\ref{n3})$,


оценка $(\ref{tri})$, условия {\em (iii), (iv)} и $(\ref{n9})$.


\end{itemize}





Тогда последовательность $(u_k)$ решений аппроксимирующих задач, построенная


выше, содержит подпоследовательность $(u_{k_l})$,


сходящуюся в ${\bold C}_1(\overline{\Omega})$ к полуправильному решению $u_0$


задачи $(\ref{uravn})$--$(\ref{gran})$,


для которого $J^{\lambda}(u_0)=\inf\limits_{v\in X}J^{\lambda}(v)$.


Если последний инфимум достигается в единственной точке $u_0$, то


$u_k\to u_0$ в ${\bold C}_1(\overline{\Omega})$.


\label{t5}


\end{thms}








\begin{pf}


Покажем, что последовательность $(u_k)$ ограничена в


${\bold C}_1(\overline{\Omega})$.


Допустим противное. Тогда существует такая подпоследовательность $(u_{k_l})$


последовательности $(u_k)$, что


$\|u_{k_l}\|_{{\bold C}_1(\overline{\Omega})}\to +\infty$. Положим


$v_l={u_{k_l}}/{\|u_{k_l}\|_{{\bold C}_1(\overline{\Omega})}}$.


Так как $u_{k_l}$ --- сильное решение задачи (\ref{uravn1})--(\ref{gran1}), то


\begin{gather}


\label{nn8}


-\sum\limits_{i,j=1}^n {(a_{ij}(x)(v_l)_{x_i})}_{x_j}=-c(x)v_l(x)+


\frac{\lambda_{k_l}g^{k_l}(x, u_{k_l}(x))}


{\|u_{k_l}\|_{{\bold C}_1(\overline{\Omega})}}, \\


%


\label{nn9}


Bv_l |_{ \Gamma}=0.


\end{gather}


Так как $\|v_l\|_{ {\bold C}_1(\overline{\Omega})}=1$, то


$\|v_l\|_{{\bold L}_q(\Omega)}\leq (\mes\Omega)^{1/q}$ $


\forall l\in {\bold N}$,


а


\begin{equation}


\label{nn10}


\frac{\|\lambda_{k_l}g^{k_l}(x, u_{k_l})\|


_{{\bold L}_q(\Omega)}}{\|u_{k_l}\|_{{\bold C}_1(\overline{\Omega})}}


\leq\frac{\lambda_{k_l}\|a\|_{{\bold L}_q(\Omega)}}


{\|u_{k_l}\|_{{\bold C}_1(\overline{\Omega})}}.


\end{equation}


Отсюда следует существование такой постоянной $M>0$, что норма правой части в


${\bold L}_q(\Omega)$ ($q>n$) равенства (\ref{nn8}) не превосходит $M$. Как


показано в (\cite{agmon}, с.\,133), последнее влечет ограниченность в 


${\bold W}_q^2(\Omega)$


последовательности $(v_l)$. Из рефлексивности ${\bold W}_q^2(\Omega)$ следует


существование подпоследовательности последовательности $(v_l)$, слабо сходящейся


к некоторому $v$ в ${\bold W}_q^2(\Omega)$. Будем обозначать ее как саму


последовательность $(v_l)$. Так как $q>n$, то ${\bold W}_q^2(\Omega)$ 


компактно вкладывается в


${\bold C}_1(\overline{\Omega})$ (\cite{sob}, с.\,103). Поэтому $v_l\to v$ в


${\bold C}_1(\overline{\Omega})$, и, значит, $Bv|_{\Gamma}=0$ (в силу


(\ref{nn9})). Из слабой сходимости $v_l$ к $v$ в ${\bold W}_q^2(\Omega)$ и оценки


(\ref{nn10}) следует


$$


\int_{\Omega}Lv(x)z(x)dx=0 \ \; \forall z\in {\bold L}_p(\Omega),


\ \ p^{-1}+q^{-1}=1.


$$


Отсюда вытекает равенство $Lv$ нулю почти всюду на $\Omega$. Таким образом,


$v$ --- ненулевое решение задачи (\ref{n6})--(\ref{n7})


($\|v\|_{{\bold C}_1(\overline{\Omega})}=1$).





Умножим обе части (\ref{uravn1}) на $v$ (при $k=k_l$ и $u=u_{k_l}$) и


проинтегрируем по $\Omega$. Получим


\begin{multline}


0=\int_{\Omega}Lu_{k_l}v(x)dx/\lambda_{k_l}=


\int_{\Omega}g^{k_l}(x, u_{k_l}(x))v(x)dx= \\


%


=\int_{\Omega}(g^{k_l}(x, u_{k_l}(x))-g(x, u_{k_l}(x)))v(x)dx+


\int_{\Omega}g(x, u_{k_l}(x))v(x)dx= \\


%


=\int_{v(x)>0}A_{k_l}(x) v(x)dx+ 


\int_{v(x)<0}A_{k_l}(x) v(x)dx+ \\


%


+\int_{v(x)>0}g(x, \|u_{k_l}\|_{{\bold C}_1


(\overline{\Omega})}v_l(x))v(x)dx+ 


%


\label{nn11}


\int_{v(x)<0}g(x, \|u_{k_l}\|_{{\bold C}_1


(\overline{\Omega})}v_l(x))v(x)dx,


\end{multline}


где $A_{k_l}|x|\le


g^{k_l}(x, \|u_{k_l}\|_{{\bold C}_1(\overline{\Omega})}v_l(x))-


g(x, \|u_{k_l}\|_{{\bold C}_1(\overline{\Omega})}v_l(x))$.





В точке $x\in\Omega$, где $v(x)>0$,


$\|u_{k_l}\|_{{\bold C}_1(\overline{\Omega})}v_l(x)\to +\infty$,


т.\,к. $v_l(x)\to v(x)$ в ${\bold C}_1(\overline{\Omega})$ и\linebreak


$\|u_{k_l}\|_{{\bold C}_1(\overline{\Omega})}\to +\infty$. Поэтому


$g^{k_l}(x, \|u_{k_l}\|_{{\bold C}_1(\overline{\Omega})}v_l(x))=


g(x, \|u_{k_l}\|_{{\bold C}_1(\overline{\Omega})}v_l(x))$


при достаточно больших $k_l$ для такого $x$ (по построению $g^{k_l}$).


Следовательно,


$$


\lim\limits_{l\to\infty}A_{k_l}(x) =0


$$


на множестве $\{x\in\Omega : v(x)>0\}$. Кроме того,


$$


A_{k_l}(x) \leq 2a(x).


$$


Поэтому по теореме Лебега о переходе к пределу под знаком интеграла получим


$$


\lim_{l\to\infty}\int_{v(x)>0}A_{k_l}(x) v(x)dx=0.


$$


Аналогично доказывается, что


$$


\lim_{l\to\infty}\int_{v(x)<0}A_{k_l}(x) v(x)dx=0.


$$





Далее, если $x$ такое, что $v(x)>0$, то 


$g(x, \|u_{k_l}\|_{{\bold C}_1(\overline{\Omega})}v_l(x))\to g_+(x)$,


т.\,к. $\|u_{k_l}\|_{{\bold C}_1(\overline{\Omega})}v_l(x)\to +\infty$


(поскольку $v_l(x)\to v(x)>0$ и 


$\|u_{k_l}\|_{{\bold C}_1(\overline{\Omega})} \to +\infty$).





Аналогично, если $v(x)<0$, то


$g(x, \|u_{k_l}\|_{{\bold C}_1(\overline{\Omega})}v_l(x))\to g_-(x)$.


Переходя к пределу в равенстве (\ref{nn11}), получим


$$


0=\int_{v(x)>0}g_+(x)v(x)dx+\int_{v(x)<0}g_-(x)v(x)dx,


$$


что противоречит условию (\ref{n9}).





Итак, ограниченность последовательности $(u_k)$ 


в ${\bold C}_1(\overline{\Omega})$ установлена.





Как отмечалось выше, $u_k$ удовлетворяет равенствам


\begin{gather}


\label{nn12}


-\sum\limits_{i,j=1}^n {(a_{ij}(x)(u_k)_{x_i})}_{x_j}=-c(x)u_k(x)+


\lambda_kg^k(x, u_k(x)),\ \ x\in\Omega, \\


%


Bu_k |_{ \Gamma}=0. \notag


\end{gather}


Так как $(u_k)$ ограничена в ${\bold C}_1(\overline{\Omega})$, то найдется


такое $M>0$, что норма правой части равенства (\ref{nn12}) в ${\bold L}_q(\Omega)$


($q>n$) не превосходит $M$ при любом $k\in {\bold N}$. Отсюда следует


ограниченность последовательности $(u_k)$ в ${\bold W}_q^2(\Omega)$ 


(\cite{agmon}, с.\,133).


Таким образом, из последовательности $(u_k)$ можно выделить слабо


сходящуюся подпоследовательность в ${\bold W}_q^2(\Omega)$ к некоторому $u_0$.


Будем ее обозначать так же, как саму последовательность $(u_k)$.


В силу компактности вложения ${\bold W}_q^2(\Omega)$ 


в ${\bold C}_1(\overline{\Omega})$ (т.\,к. $q>n$)


получаем сильную сходимость $(u_k)$ к $u_0$ в ${\bold C}_1(\overline{\Omega})$.


Отсюда, в частности, следует равенство $Bu_0\vert{}_{\Gamma}=0$.





Докажем, что


$J^\lambda(u_0)=\inf\limits_{v\in X}J^\lambda(v)=d_\lambda$.


В силу теоремы \ref{t3} существует $\wh{u}_0\in X$ такое, что


$J^\lambda(\wh{u}_0)=d_\lambda$.





Заметим, что для любого натурального $l$


\begin{multline*}


|J_l(u_l)-J^\lambda(u_l)|=


\bigg|\lambda_l\int\limits_{\Omega}dx\int\limits_0^{u_l(x)}g^l(x, s)ds-


\lambda\int\limits_{\Omega}dx\int\limits_0^{u_l(x)}g(x, s)ds\bigg|\leq \\


%


\le |\lambda_l-\lambda|\int\limits_{\Omega}|u_l(x)|a(x)dx+|\lambda|


\bigg|\int\limits_{\Omega}dx\int\limits_0^{u_l(x)}


(g^l(x, s)-g(x, s))ds\bigg|\leq \\


% 


\le |\lambda_l-\lambda|\,\|u_l\|_{{\bold L}_p(\Omega)}


\|a\|_{{\bold L}_q(\Omega)}+ 


|\lambda|\sum_{i=1}^m\int\limits_{\Omega}dx


\int\limits_{\varphi_i(x)-\delta_l}^{\varphi_i(x)+\delta_l}


|g^l(x, s)-g(x, s)|ds\leq \\


%


\le |\lambda_l-\lambda|\,\|u_l\|_{{\bold L}_p(\Omega)}\|a\|_{{\bold L}_q(\Omega)}+


4|\lambda|m\delta_l\int\limits_{\Omega}a(x)dx=\wh{\gamma_l},\quad p^{-1}+q^{-1}=1.


\end{multline*}


Аналогично доказывается, что


$$


|J_l(\wh{u}_0)-J^\lambda(\wh{u}_0)|\leq |\lambda_l-\lambda|\,


\|\wh{u}_0\|_{{\bold L}_p(\Omega)}\|a\|_{{\bold L}_q(\Omega)}


+4|\lambda|m\delta_l\int_{\Omega}a(x)dx=\wh{\wh{\gamma_l}}.


$$





Положим $\gamma_l=\max\{\wh{\gamma_l}, \wh{\wh{\gamma_l}}\}$.


Так как $(u_l)$ ограничена в ${\bold C}_1(\overline{\Omega})$, то


$\gamma_l\to 0+$.





Покажем, что $J_l(u_l)\to d_\lambda$. В силу оценок


$|J_l(u_l)-J^\lambda(u_l)|\leq\gamma_l$ и


$|J_l(\wh{u}_0)-J^\lambda(\wh{u}_0)|\leq\gamma_l$


имеем


$d_\lambda-\gamma_l\leq J^\lambda(u_l)-\gamma_l\leq J_l(u_l)\leq


J_l(\wh{u}_0)\leq J^\lambda(\wh{u}_0)+\gamma_l=d_\lambda+\gamma_l$.


Таким образом, $d_\lambda-\gamma_l\leq J_l(u_l)\leq d_\lambda+\gamma_l$, что 


влечет сходимость $J_l(u_l)$ к $d_\lambda$.





Функционал $J^\lambda$ непрерывeн на $X$, поэтому


$\lim\limits_{l\to\infty}J^\lambda(u_l)=J^\lambda(u_0)$,


поскольку $u_l\to u_0$ в $X$. С другой стороны,


$\lim\limits_{l\to\infty}J^\lambda(u_l)=


\lim\limits_{l\to\infty}J_l(u_l)=d_\lambda$.


Таким образом, $J^\lambda(u_0)=d_\lambda$.





Согласно теореме \ref{t3} \ $u_0$ является полуправильным решением задачи


(\ref{uravn})--(\ref{gran}). Из приведенного выше доказательства видно, что из


любой подпоследовательности последовательности $(u_k)$ сильных решений


аппроксимирующих задач (\ref{uravn1})--(\ref{gran1}) можно выделить


подпоследовательность, сходящуюся в ${\bold C}_1(\overline{\Omega})$ к некоторому


$u_0\in X$, для которого верно (\ref{n8}). Если элемент $u_0$, удовлетворяющий


(\ref{n8}), единственный, то отсюда стандартным рассуждением получаем, что сама


$(u_k)$ сходится к $u_0$ в ${\bold C}_1(\overline{\Omega})$.


\end{pf}








\begin{thebibliography}{9}





\bibitem{krasn2}


Красносельский М.А., Покровский А.В. {\it Уравнения с разрывными


нелинейностями} // ДАН СССР. -- 1979. -- Т.\,248. -- \No\,5. -- C.\,1056--1059.


\bibitem{pavl4}


Павленко В.Н., Искаков Р.С. {\it Непрерывные аппроксимации разрывных


нелинейностей полулинейных уравнений эллиптического типа} // Укр. матем. журн.


-- 1999. -- Т.\,51. -- \No\,2. -- С.\,224--233.


\bibitem{pot1}          


Павленко В.Н., Потапов Д.К. {\it О существовании луча собственных значений


для уравнений с разрывными операторами} // Сиб. матем. журн. -- 2001. -- Т.\,42.


-- \No\,4. -- С.\,911--919.


\bibitem{lad}


Ладыженская О.А., Уральцева Н.Н. {\it Линейные и квазилинейные уравнения


эллиптического типа.} -- М.: Наука, 1964. -- 538~с.


\bibitem{sob}


Соболев С.Л. {\it Некоторые применения функционального анализа в


математической физике.} -- М.: Наука, 1982. -- 336~с.


\bibitem{kar}


Красносельский М.А., Покровский А.В. {\it Системы с гистерезисом.} --


М.: Наука, 1983. -- 272~с.


\bibitem{krasn3}


Красносельский М.А., Покровский А.В. {\it Правильные решения уравнений с


разрывными нелинейностями} // ДАН СССР. -- 1976. -- Т.\,226. -- \No\,3. --


C.\,506--509.


\bibitem{vadim}


Павленко В.Н., Винокур В.В. {\it Резонансные краевые задачи для уравнений


эллиптического типа с разрывными нелинейностями} // Изв. вузов. Математика. 


-- 2001. -- \No\,5. -- С.\,43--58.


\bibitem{agmon}


Агмон С., Дуглис А., Ниренберг Л. {\it Оценки решений эллиптических уравнений


вблизи границы.} -- М.: Ин. лит., 1962. -- 208~с.





\end{thebibliography}





\vskip 1cm





\post{\it Челябинский государственный}{\it Поступила}


\post{\it университет}{26.02.2003}





\label{end}


\end{document}


